%\noindent\resizebox{\textwidth}{!}{%
%\begin{tikzpicture}
%\tikzset{
%  hdr/.style={fill=black!15, font=\footnotesize\bfseries, minimum height=0.7cm, text centered, inner sep=3pt},
%  cel/.style={font=\footnotesize, minimum height=1.0cm, text centered, inner sep=3pt, text width=#1, align=center},
%}
%
%% Column widths
%\def\colA{4.6cm}
%\def\colB{3.8cm}
%\def\colC{3.2cm}
%\def\colD{3.6cm}
%
%% Header row
%\node[hdr, text width=\colA] (ha) at (0,0) {Блок моделі};
%\node[hdr, text width=\colB, right=-\pgflinewidth of ha] (hb) {Фактори CAMIL};
%\node[hdr, text width=\colC, right=-\pgflinewidth of hb] (hc) {Компоненти TPACK};
%\node[hdr, text width=\colD, right=-\pgflinewidth of hc] (hd) {Етап 6-комп. моделі};
%
%% Row 1: Мотиваційно-цільовий
%\node[cel=\colA, below=-\pgflinewidth of ha] (a1) {Мотиваційно-цільовий};
%\node[cel=\colB, right=-\pgflinewidth of a1] (b1) {мотивація,\\саморегуляція};
%\node[cel=\colC, right=-\pgflinewidth of b1] (c1) {PK, CK};
%\node[cel=\colD, right=-\pgflinewidth of c1] (d1) {Аналіз};
%
%% Row 2: Теоретико-методологічний
%\node[cel=\colA, below=-\pgflinewidth of a1] (a2) {Теоретико-\\методологічний};
%\node[cel=\colB, right=-\pgflinewidth of a2] (b2) {когнітивне\\навантаження};
%\node[cel=\colC, right=-\pgflinewidth of b2] (c2) {TK+PK+CK};
%\node[cel=\colD, right=-\pgflinewidth of c2] (d2) {Вимоги};
%
%% Row 3: Змістовий
%\node[cel=\colA, below=-\pgflinewidth of a2] (a3) {Змістовий};
%\node[cel=\colB, right=-\pgflinewidth of a3] (b3) {втілення,\\агентність};
%\node[cel=\colC, right=-\pgflinewidth of b3] (c3) {CK, TK};
%\node[cel=\colD, right=-\pgflinewidth of c3] (d3) {Підбір ПЗ +\\планування};
%
%% Row 4: Організаційно-технологічний
%\node[cel=\colA, below=-\pgflinewidth of a3] (a4) {Організаційно-\\технологічний};
%\node[cel=\colB, right=-\pgflinewidth of a4] (b4) {агентність,\\присутність};
%\node[cel=\colC, right=-\pgflinewidth of b4] (c4) {TK, TPK};
%\node[cel=\colD, right=-\pgflinewidth of c4] (d4) {Розробка +\\впровадження};
%
%% Row 5: Рефлексивно-оцінний
%\node[cel=\colA, below=-\pgflinewidth of a4] (a5) {Рефлексивно-оцінний};
%\node[cel=\colB, right=-\pgflinewidth of a5] (b5) {саморегуляція};
%\node[cel=\colC, right=-\pgflinewidth of b5] (c5) {TPK, TPACK};
%\node[cel=\colD, right=-\pgflinewidth of c5] (d5) {Оцінювання};
%
%% Draw grid borders
%\foreach \r in {a,b,c,d} {
%  \foreach \i in {1,2,3,4,5} {
%    \draw (\r\i.north west) rectangle (\r\i.south east);
%  }
%}
%% Header borders
%\draw (ha.north west) rectangle (ha.south east);
%\draw (hb.north west) rectangle (hb.south east);
%\draw (hc.north west) rectangle (hc.south east);
%\draw (hd.north west) rectangle (hd.south east);
%
%% Outer frame
%\draw[thick] (ha.north west) rectangle (d5.south east);
%% Thick line below header
%\draw[thick] (ha.south west) -- (hd.south east);
%
%\end{tikzpicture}%
%}
{\small
	\begin{tabularx}{\linewidth}{|p{4.3cm} |l |c|X|}
		\hline
		\rowcolor{gray!20}
		\hfil \textbf{Блок моделі} & 
		\hfill \textbf{Фактори CAMIL} & 
		\textbf{\makecell{Компоненти \\TPACK}} & 
		\hfil \textbf{Етап 6-комп. моделі} \\
		\hline
		
		Мотиваційно-цільовий 
		& мотивація, саморегуляція 
		& PK, CK 
		& Аналіз \\
		\hline
		Теоретико-методологічний 
		& когнітивне навантаження 
		& TK+PK+CK 
		& Вимоги \\
		\hline
		Змістовий 
		& втілення, agentність 
		& CK, TK 
		& Підбір ПЗ + планування \\
		\hline
		Організаційно-технологічний 
		& агентність, присутність 
		& TK, TPK 
		& Розробка + впровадження \\
		\hline
		Рефлексивно-оцінний 
		& саморегуляція 
		& TPK, TPACK 
		& Оцінювання \\
		\hline
	\end{tabularx}
}